\chapter{Vladimir G. Drinfeld}\label{ch1990vladimirdrinfeld:vladimir-g.-drinfeld-fields-medal-1990-a-comprehensive-report}

\begin{center}\rule{0.5\linewidth}{0.5pt}\end{center}

\textbf{Prepared:} July 2026 \textbf{Subject:} Vladimir Gershonovich Drinfeld (b.~February 14, 1954, Kharkiv, Ukraine, USSR) \textbf{Position:} Professor, University of Chicago (since 1981) \textbf{Doctoral Advisor:} Yuri Manin (Steklov Institute) \textbf{Fields Medal Citation:} \emph{``For his work on the Langlands program, particularly his proof of the geometric Langlands conjecture for function fields.''}

\begin{center}\rule{0.5\linewidth}{0.5pt}\end{center}

\begin{prereqs} Algebraic geometry, number theory (class field theory, automorphic forms), representation theory, Langlands program, algebraic groups. \end{prereqs}

\begin{objectives} State the geometric Langlands conjecture for function fields; explain the role of Drinfeld modules in function field arithmetic; describe the connection between the geometric Langlands program and the classical Langlands program. \end{objectives}

\section{Problem History}

The Langlands program, formulated by Robert Langlands in the 1960s, is a vast set of conjectures connecting number theory, representation theory, and algebraic geometry. The central conjecture---the automorphic-to-Galois direction of the Langlands correspondence---asserts that automorphic representations of $GL_n$ over a number field correspond to $n$-dimensional Galois representations.

The geometric Langlands conjecture is the function field analogue of the Langlands program. For a function field $F = \mathbb{F}_q(C)$ (the field of rational functions on a curve $C$ over a finite field $\mathbb{F}_q$), the geometric Langlands correspondence is a relationship between automorphic sheaves on the moduli of $GL_n$-bundles and local systems (Galois representations) on the curve.

The arithmetic Langlands program is the number field version, where one works over number fields like $\mathbb{Q}$. The geometric version is ``geometric'' because it involves sheaves and local systems on algebraic varieties, rather than representations of number-theoretic Galois groups. The geometric Langlands conjecture was formulated by Drinfeld and Laumon \cite{drinfeld1987} in the 1980s.

Before Drinfeld's work, the geometric Langlands conjecture was known only in special cases. For $GL_2$, the conjecture was proved by Drinfeld in 1974 for function fields, using the theory of Drinfeld modules. For general $n$, the conjecture was proved by Laumon, Drinfeld, and others in the 1990s.

Drinfeld modules, introduced by Drinfeld in 1974, are the function field analogue of elliptic curves. A Drinfeld module of rank $r$ over a function field $F$ is a ring homomorphism $\phi: \mathbb{F}_q[T] \to \text{End}_F(\mathbb{G}_a)$ from the polynomial ring to the endomorphism ring of the additive group, with certain properties. The rank $r$ is the dimension of the $F$-vector space $\phi(T) \cdot F$.

Drinfeld modules provide the function field analogue of modular forms and elliptic curves. The theory of complex multiplication for Drinfeld modules---the analogue of the theory of complex multiplication for elliptic curves---is the foundation for Drinfeld's proof of the Langlands correspondence for $GL_2$ over function fields.

The geometric Langlands conjecture, in its modern form, states: for a reductive group $G$ over a function field $F$, the category of $G$-local systems on the curve $C$ is equivalent to the category of $G$-automorphic sheaves on the moduli stack of $G$-bundles on $C$. This is a categorical version of the Langlands correspondence, involving sheaves and perverse sheaves rather than sets of representations.

The proof of the geometric Langlands conjecture uses the theory of perverse sheaves, the geometric Satake equivalence, and the geometric Langlands correspondence as developed by Drinfeld and Laumon \cite{drinfeld1987}. The key idea is to relate the geometric Langlands correspondence to the classical Langlands correspondence through the theory of automorphic sheaves.

The historical context of Drinfeld's work includes the earlier work of Eichler and Shimura on the correspondence between modular forms and elliptic curves, and the work of Drinfeld's advisor Manin on automorphic forms and algebraic geometry. Drinfeld's innovation was to replace number fields by function fields, where the geometry of curves over finite fields provides a more accessible setting.

\section{How Was It Solved}

Drinfeld's proof of the Langlands correspondence for $GL_2$ over function fields proceeds in three steps: (1) construct automorphic forms on $GL_2$ over function fields; (2) associate to each automorphic form a Galois representation; (3) show the correspondence is a bijection.

The automorphic forms are functions on the double coset $GL_2(F) \backslash GL_2(\mathbb{A}_F) / GL_2(\hat{\mathcal{O}})$, where $\mathbb{A}_F$ is the adele ring of $F$ and $\hat{\mathcal{O}}$ is the product of the rings of integers at all places. The Galois representations are 2-dimensional representations of $\text{Gal}(\overline{F}/F)$.

The key construction is the use of Drinfeld modules. Given a Drinfeld module $\phi$ of rank 2, one associates a 2-dimensional Galois representation $\rho_\phi: \text{Gal}(\overline{F}/F) \to GL_2(\mathbb{C})$. The construction uses the theory of torsion points of $\phi$: the $\mathfrak{p}$-torsion $\phi[\mathfrak{p}]$ (for a prime $\mathfrak{p}$ of $\mathbb{F}_q[T]$) is a 2-dimensional $\mathbb{F}_\mathfrak{p}$-vector space, and the Galois group acts on these torsion points.

The geometric Langlands conjecture, proved by Drinfeld and Laumon \cite{drinfeld1987}, states: the category of $GL_n$-local systems on $C$ is equivalent to the category of $GL_n$-automorphic sheaves on the moduli of $GL_n$-bundles. The proof uses the geometric Satake equivalence, which identifies the category of spherical perverse sheaves on the affine Grassmannian with the category of representations of the Langlands dual group.

The geometric Satake equivalence, proved by Ginzburg, Mirkovic, and Vilonen, is the geometric analogue of the classical Satake isomorphism. It states that the category of spherical perverse sheaves on the affine Grassmannian $\text{Gr}_G$ is equivalent to the category of representations of the Langlands dual group ${}^LG$. This equivalence is used to transfer between automorphic sheaves and local systems.

The proof of the geometric Langlands correspondence uses the theory of Hecke operators. The Hecke operators are morphisms between the moduli spaces of $G$-bundles, and they generate a commutative algebra acting on the cohomology of these moduli spaces. The eigenspaces of the Hecke operators correspond to local systems, giving the Langlands correspondence.

Drinfeld also introduced the notion of a Drinfeld module as a tool for studying the arithmetic of function fields. The theory of Drinfeld modules generalizes the theory of elliptic curves to function fields, providing function field analogues of many results in number theory: class field theory, the theory of complex multiplication, and the Langlands program.

The key property of Drinfeld modules that makes them useful for the Langlands correspondence is the existence of a canonical ring homomorphism:
\[
\rho_\phi: \text{Gal}(\overline{F}/F) \to GL_r(\hat{\mathbb{F}}_q[[T]])
\]
where $\hat{\mathbb{F}}_q[[T]]$ is the completion of $\mathbb{F}_q[T]$ at the maximal ideal $(T)$. This Galois representation is compatible with the action of the Hecke operators at all places, providing the automorphic structure.

The geometric Langlands correspondence has applications to the theory of automorphic forms and the study of the Langlands program in general. The geometric version provides a geometric framework for the Langlands program that is often more accessible than the arithmetic version.

\section{Impact of This Discovery}

Drinfeld's work on the Langlands program opened a new chapter in the relationship between number theory and algebraic geometry. The geometric Langlands correspondence provides a bridge between the theory of sheaves and the theory of automorphic forms, with applications to both fields.

The theory of Drinfeld modules has become a fundamental tool in function field arithmetic. The function field analogue of Fermat's last theorem, proved by Drinfeld and later by others, uses Drinfeld modules as a key ingredient. The theory of Drinfeld modules also has applications to the theory of modular forms and the study of $L$-functions.

The geometric Langlands program has led to many new developments in algebraic geometry. The theory of perverse sheaves, the geometric Satake equivalence, and the study of the moduli of $G$-bundles are all central topics in modern algebraic geometry. The geometric Langlands correspondence provides a unifying framework for many of these developments.

The relationship between the geometric Langlands program and the arithmetic Langlands program is an active area of research. The geometric version often provides insights and tools for the arithmetic version. For example, the proof of the function field Langlands conjecture for $GL_n$ uses ideas from the geometric theory.

Drinfeld's work has also influenced the theory of quantum groups and vertex algebras. The theory of Drinfeld modules is closely related to the theory of quantum groups, and the geometric Langlands correspondence has applications to the study of conformal field theory and vertex operator algebras.

The Langlands program remains one of the most important and active areas of research in mathematics. Drinfeld's contributions---both the proof of the Langlands correspondence for function fields and the development of the theory of Drinfeld modules---have been fundamental to the progress in this field.

The theory of shtukas, introduced by Drinfeld, is the function field analogue of the Langlands correspondence. A shtuka of rank $n$ on a curve $C$ is a pair $(\mathcal{F}_1, \mathcal{F}_2)$ of $GL_n$-bundles together with an isomorphism away from finitely many points. The moduli space of shtukas is the function field analogue of the modular curve, and its cohomology realizes the Langlands correspondence.

The function field analogue of the Ramanujan conjecture, proved by Drinfeld using Drinfeld modules, states that the Hecke eigenvalues of automorphic forms on $GL_2$ over function fields satisfy the Ramanujan bound. This result, proved before the arithmetic analogue by Deligne, was one of the early triumphs of the function field approach. Specifically, if $f$ is a normalized Hecke eigenform of weight $k$ on $\text{GL}_2(\mathbb{A}_F)$, then the eigenvalue $\lambda_\mathfrak{p}$ at a place $\mathfrak{p}$ satisfies $|\lambda_\mathfrak{p}| \leq 2q_\mathfrak{p}^{(k-1)/2}$, where $q_\mathfrak{p}$ is the size of the residue field.

Drinfeld's proof of the Ramanujan conjecture uses the theory of Drinfeld modules as a tool for constructing Galois representations with prescribed local properties. The key idea is that the Hecke eigenvalues appear as traces of Frobenius elements in the Galois representation associated to a Drinfeld module, and the Ramanujan bound follows from the Weil conjectures applied to the moduli space of Drinfeld modules.

\section{Modern Views and Improvements}

The geometric Langlands correspondence has been extended to many other settings. The affine and mixed versions of the geometric Langlands program, studied by many researchers, provide a more complete picture of the correspondence. The relationship to the theory of automorphic forms on $p$-adic groups, through the local Langlands correspondence, is an active area of research.

The theory of Drinfeld modules has been generalized in many directions. Drinfeld modules of higher rank, shtukas (the function field analogue of Galois representations), and the function field analogue of the Iwasawa theory have all been developed. The theory of shtukas, introduced by Drinfeld, provides a geometric framework for the Langlands program.

The geometric Langlands correspondence has applications to the theory of mirror symmetry. The Langlands dual group ${}^LG$ appears naturally in the description of the mirror partner of a variety related to the moduli of $G$-bundles, and the geometric Langlands equivalence can be interpreted as a form of mirror symmetry. This connection, developed by Drinfeld, Beilinson, and others, provides a bridge between the Langlands program and string theory.

The theory of vertex algebras, which Drinfeld has studied extensively, provides a mathematical framework for conformal field theory. Vertex algebras are algebraic structures that encode the operator product expansion of chiral fields in 2-dimensional CFT, and they have applications to the theory of moonshine, the study of the monster group, and the construction of new invariants of knots and 3-manifolds. Drinfeld's work on vertex algebras has influenced the development of the geometric Langlands program and the theory of quantum groups.

The relationship between the geometric Langlands program and the theory of D-modules is fundamental. The category of D-modules on the moduli of $G$-bundles provides a natural home for the geometric Langlands correspondence, and the study of D-modules on singular spaces has led to new insights into the structure of the correspondence. The theory of D-modules, as developed by Bernstein, Beilinson, and others, provides the technical tools for the geometric Langlands program.

The arithmetic Langlands program has seen major advances in recent years. The proof of the Sato--Tate conjecture by Taylor, Harris, and Shepherd-Barron, and the proof of the modularity theorem by Wiles, Taylor, and others, are major milestones in the Langlands program. The geometric Langlands program provides a framework for understanding these results.

Drinfeld himself has continued to make major contributions to mathematics. His work on vertex algebras, the geometric Langlands program, and the theory of automorphic forms has influenced many areas of mathematics and mathematical physics. The Langlands program remains one of the most important and active areas of research in mathematics.

The function field approach to the Langlands program continues to be a source of new ideas. The geometric Langlands correspondence, as formulated by Drinfeld and Laumon \cite{drinfeld1987}, has been generalized to include wild ramification, mixed characteristics, and higher-dimensional base spaces. The relationship between the geometric Langlands program and the theory of factorization algebras, as developed by Beilinson and Drinfeld, provides a categorical framework for the correspondence.

The Drinfeld upper half-plane $\Omega^n = \mathbb{P}^{n-1} \setminus \bigcup_{H \subset \mathbb{F}_q^n} H$, the function field analogue of the upper half-plane, is a fundamental object in Drinfeld's work. The cohomology of the Drinfeld upper half-plane realizes the Langlands correspondence, and its study connects to the theory of modular forms, automorphic representations, and the Langlands program.

\begin{center}\rule{0.5\linewidth}{0.5pt}\end{center}

\section*{Summary}
\addcontentsline{toc}{section}{Summary}

Drinfeld proved the Langlands correspondence for function fields using Drinfeld modules and introduced the geometric Langlands conjecture, connecting local systems to automorphic sheaves via the geometric Satake equivalence. His work on Drinfeld modules established the function field analogue of elliptic curves and modular forms, opening new directions in arithmetic geometry and the Langlands program.
